\chapter*{Abstract}
\addcontentsline{toc}{chapter}{Abstract}

Graph Neural Operators (GNOs) can construct smooth implied volatility surfaces from irregular option quotes in a single forward pass, but a deterministic GNO does not quantify uncertainty around its estimate. This thesis develops a post-hoc conformal prediction framework that adds calibrated lower and upper surfaces to a fixed, pre-trained GNO without modifying or retraining the model. Two economically motivated nonconformity scores are compared: a spread-normalised price residual (Score A) and a Vega-weighted implied volatility residual (Score B). Calibration uses the five most recent market surfaces, while a $4\times4$ Mondrian partition assigns spatially varying quantiles across maturity and normalised moneyness. The bands are subsequently adjusted together to enforce a calendar condition and a discrete butterfly condition based on neighbouring grid points, while keeping implied volatility positive and ensuring that the lower bound remains below the upper bound.

The empirical study uses 46 retained S\&P 500 option surfaces collected between December 2025 and August 2026. The first 15 surfaces form the initial calibration period and the remaining 31 are evaluated in walk-forward order. On each surface, 80\% of the quotes are supplied to the GNO as context and 20\% are withheld for calibration or evaluation. With a nominal coverage level of 90\%, Score B achieves mean post-projection coverage of 90.80\% under global calibration and 92.85\% under Mondrian calibration, compared with 85.46\% and 84.27\%, respectively, for Score A. The mean band widths are 0.027180 implied-volatility units for Score B and 0.067806 for Score A, making the Score A bands approximately 2.49 times wider. The projection reduces all violations of the two enforced grid conditions to zero for both scores and produces no crossings at the evaluated target quotes.

Within this dataset, Score B provides the stronger combination of empirical coverage, sharpness, and stability across surfaces. The results show that conformal prediction can serve as a practical uncertainty layer for neural-operator volatility smoothing. Because the observations are temporally ordered, the reported coverage is an empirical walk-forward result rather than a claim of the standard exchangeable-data guarantee.
